% !TEX program = pdflatex
\documentclass[12pt,a4paper]{letter}

\usepackage[T1]{fontenc}
\usepackage[margin=1in]{geometry}
\usepackage{hyperref}
\usepackage{setspace}
\usepackage{parskip}

\hypersetup{
  colorlinks=true,
  urlcolor=blue
}

\onehalfspacing

% Push closing/signature block to the right margin
\longindentation=0.65\textwidth

\signature{Professor Ayo Ogunjuyigbe\\
Professor of Electrical \& Electronic Engineering\\
Provost, Postgraduate School\\
University of Ibadan, Nigeria}

\address{Postgraduate School\\
University of Ibadan\\
Ibadan, Oyo State\\
Nigeria}

\date{27 February 2026}

\begin{document}

\begin{letter}{UK Visas \& Immigration\\
Global Talent Visa --- Tech Nation (Digital Technology)\\
United Kingdom}

\opening{To Whom It May Concern,}

I am writing to endorse Odeyemi Olusegun Israel for the UK Global Talent Visa under the Digital Technology route. I have mentored Olusegun for a number of years and have followed his technical work closely---particularly so over the last twelve months, during which the scope of that work expanded into areas I know well.

I am a Professor of Electrical and Electronic Engineering at the University of Ibadan, where I currently serve as Provost of the Postgraduate School. My research is in power systems, computational intelligence, and applied machine learning for engineering. That background matters here, because Olusegun's most recent paper includes results on power-grid instability that fall squarely within my area of expertise, and I want to comment on them directly.

\textbf{The AMTTP Platform.} Olusegun has designed and built, largely on his own, a system called the Anti-Money Laundering Transaction Trust Protocol (AMTTP). It is a four-layer architecture: client SDKs (TypeScript and Python, both open-sourced), a compliance orchestration engine, an ML training pipeline, and a blockchain smart contract layer with 18 contracts deployed on Ethereum Sepolia. Seventeen containerised microservices sit behind a gateway. The system tackles a problem I consider genuinely unsolved in decentralised finance: enforcing regulatory compliance \textit{before} a transaction settles, not after.

Every major AML tool on the market---Chainalysis, Elliptic, and their competitors---works post-settlement. They flag; they do not prevent. AMTTP intervenes at the protocol level, combining ML risk scoring, graph analysis, sanctions screening, and policy adjudication into a single deterministic pass/fail decision. The training pipeline processes over 2.6 million transactions through a Composite Teacher model blending Autoencoder-Enhanced XGBoost, FATF-aligned pattern detectors, and graph structural features. He also built a zero-knowledge proof framework (zkNAF) for privacy-preserving credential verification, along with multi-oracle threshold signatures and replay protection.

I do not say any of this lightly. I have seen plenty of systems that stitch together existing libraries and call it innovation. AMTTP is different. Someone sat with the regulatory text, understood the engineering constraints of blockchain finality, and designed a system from scratch to reconcile the two. That takes a particular kind of patience and clarity.

\textbf{Original Research Contributions.} The earlier papers are solid work on their own. His \textit{Blind Spot Decomposition Theory} (BSDT)\footnote{Available at \url{https://doi.org/10.5281/zenodo.18870589}.} decomposes the failure modes of ML detectors into four near-orthogonal components---Camouflage, Feature Gap, Activity Anomaly, and Temporal Novelty---accounting for over 95\% of explainable variance in missed-fraud indicators. The correction operators recover recall by up to 84.4 percentage points without retraining, across XGBoost, LightGBM, Random Forest, and Logistic Regression on six datasets in five domains. His \textit{Universal Detection Principle} (UDP)\footnote{Available at \url{https://doi.org/10.5281/zenodo.19037687}.} is a tensor-based anomaly detection framework that complements BSDT: where BSDT explains \textit{why} a model fails, UDP measures deviation across multiple representations simultaneously to catch what single-spectrum approaches miss.

But it is his latest paper---\textit{Blind-Spot Decomposition and the Geometry of System Collapse}\footnote{Available at \url{https://doi.org/10.5281/zenodo.19038783}.}---that changed how I think about his work.

\textbf{The geometry paper and the power-grid results.} In this paper, Olusegun proves that the BSDT energy functional is topologically equivalent to the interaction potential of a coupled dynamical system. The core result, which he calls Theorem~C, establishes Morse-index equivalence at every critical point, homotopy equivalence of sublevel sets, and two-sided gradient bounds. The practical consequence is significant: detecting instability and controlling it turn out to be two views of the same geometric object. He then derives a spectral phase-transition theorem showing that below a critical manifold, any constant coupling stabilises the system, but above it, no fixed parameter works---only adaptive friction succeeds. If you are a power engineer, the implications of that phase-transition result are hard to ignore.

Here is what caught my attention most. He applies this framework---built originally for financial networks---to ERCOT power-grid dispatch data and achieves an AUC of 0.9997 with a 72-hour lead time before the rolling blackouts of Winter Storm Uri. The adaptive friction controller reduces peak capacity loss by 71\%. There is no retraining, no domain-specific tuning. The same four BSDT channels that detect financial fraud map directly onto power-grid instability: enstrophy maps to camouflage, spectral cascade to feature gap, vorticity-strain alignment to activity anomaly, temporal acceleration to temporal novelty. I have worked in power systems for over two decades. A 72-hour lead time with near-perfect discrimination and zero domain-specific adjustment is not something I have seen before.

He goes further. The same framework is applied to three-dimensional incompressible Navier--Stokes equations, validated through a 34-run GPU stress test on an NVIDIA RTX PRO 6000 Blackwell at resolutions up to $256^3$ and Reynolds numbers up to 62,832. All 34 runs complete without blow-up. Adaptive viscosity reduces peak vorticity by 1.77$\times$ at $\mathrm{Re} \approx 1257$; the BKM integral stays finite at $\mathrm{Re} \approx 62{,}832$; and five distinct feedback laws produce dynamics identical to within 0.04\% spread. He is careful not to claim unconditional regularity---the results are conditional and numerical---but the fact that a framework designed for financial networks transfers to fluid dynamics without modification is, in my view, a striking result.

The paper also validates on FDIC banking data (6-quarter lead on the Global Financial Crisis, AUROC of 0.867, zero false alarms) and the Terra/Luna crypto collapse ($r = 0.996$ correlation with historical prices, 30-hour advance warning). The breadth is unusual. Nine contributions, three entirely separate physical domains, and a single underlying geometric structure connecting them. I have reviewed a good number of journal manuscripts, and I cannot recall the last time I saw this kind of range from a single author.

\textbf{My assessment.} I have taught and mentored hundreds of engineering students and researchers at Ibadan. A handful, over the years, have genuinely surprised me. Olusegun is one of them.

His earlier work on AMTTP and the original BSDT paper would have been enough for a strong endorsement. But the geometry paper moves him into different territory. He has produced a piece of applied mathematics that proves a geometric equivalence between detection and control in coupled systems, validates it on real power-grid data with results I find credible and impressive, and then extends it to Navier--Stokes. He did this without a university affiliation, without co-authors, and without grant funding. I am not sure what else to call that except exceptional.

The UK has a direct interest here. The FCA's forthcoming crypto framework (expected 2027) will need people who understand compliance engineering at the protocol level---Olusegun already builds that. The National Grid's decarbonisation and resilience challenges will need people who can bring rigorous mathematical tools to bear on power-system instability---Olusegun's ERCOT results show he can do that too. He is not a specialist in one narrow area waiting for the right problem. He is someone who builds frameworks and then proves they work in domains he was not originally trained in. That is the sort of person you want.

I endorse his application with confidence.

\vspace{1.5\baselineskip}

\closing{Yours faithfully,}

\end{letter}

\end{document}
